% Generated by Listas/make_statements.py from pratica-lecture-01.md, item 13. Do not edit.
\textbf{Source.} Mas-Colell, Whinston \& Green (1995), \emph{Microeconomic Theory}, Ch. 8 Exercises, Exercise 8.B.1\^{}A, p.~262 {[}PDF 142{]}. Reproduced verbatim. Answer in the source: none in the book (see the MWG solutions manual in \texttt{Livros/}).



\begin{quote}

There are \(I\) firms in an industry. Each can try to convince Congress to give the industry a subsidy. Let \(h_i\) denote the number of hours of effort put in by firm \(i\), and let \(c_i(h_i) = w_i(h_i)^2\), where \(w_i\) is a positive constant, be the cost of this effort to firm \(i\). When the effort levels of the firms are \((h_1, \ldots, h_I)\), the value of the subsidy that gets approved is \(\alpha \sum_i h_i + \beta(\prod_i h_i)\), where \(\alpha\) and \(\beta\) are constants.



Consider a game in which the firms decide simultaneously and independently how many hours they will each devote to this effort. Show that each firm has a strictly dominant strategy if and only if \(\beta = 0\). What is firm \(i\)'s strictly dominant strategy when this is so?

\end{quote}

